% part: intuitionistic-logic
% chapter: propositions-as-types

\documentclass[../../../include/open-logic-chapter]{subfiles}

\begin{document}

\olchapter{int}{pty}{Proposisi minangka Tipe}

\begin{editorial}
  Iki dhraf bab ngenani korespondhensi Curry--Howard sing
  \emph{eksperimental banget}. Dhraf iki isih mbutuhake panjelasan
  lan motivasi liyane, lan bisa uga ana kaluputan lan bagean sing
  kurang. Bukti normalisasi kudu ditliti lan dijangkepi. Ora ana
  conto kanggo tipe prodhuk. Konversi permutasi lan panyederhanaan
  durung dibahas. Bab iki bakal luwih cetha yen uga ana materi ngenani
  kalkulus lambda (mawa tipe), sing ing kene dianggep wis dimangerteni.
  Gunakna kanthi ngati-ati banget.
\end{editorial}

\olimport{introduction}
%\olimport{sequent-natural-deduction}
\olimport{proof-terms}
\olimport{proofs-to-terms}
\olimport{terms-to-proofs}
\olimport{types}
\olimport{reduction}
\olimport{type-preservation}
\olimport{normalization}

\OLEndChapterHook

\end{document}
